%%%%%%%%%%%%%%%%%%%%%%%%%%%%%%%%%%%%%%%%%%%%%%%%%%%%%%%%%%%%%%%%%%%%%%
% Lớp 11 | Một số yếu tố vẽ kỹ thuật | Hình biểu diễn của một hình, khối
% Dạng 1C3?1-2: Lý thuyết về phép chiếu vuông góc
% Số câu: 12  (sinh tự động bởi tach_cau.py; từ nay thầy sửa trực tiếp file này)
%%%%%%%%%%%%%%%%%%%%%%%%%%%%%%%%%%%%%%%%%%%%%%%%%%%%%%%%%%%%%%%%%%%%%%

%%%%%%%%%%%%%%%% TỰ LUẬN %%%%%%%%%%%%%%%%

% nguồn: Ok_TuanLuan_Toan11_Lan1
\begin{bt}%[1C3Y1-2]%[ma:0003c034]
	 Tìm giới hạn sau $\lim\dfrac{3^{n+1}-4^n}{4^{n-1}+3}$.
	\loigiai{ 
	\[\lim\dfrac{3^{n+1}-4^n}{4^{n-1}+3} = \lim\dfrac{9\cdot 3^{n-1}-4\cdot  4^{n-1}}{4^{n-1}+3}=\lim\dfrac{9\cdot\left(\dfrac{3}{4}\right)^{n-1}-4}{1+3\cdot\left(\dfrac{1}{4}\right)^{n-1}}=-4.\]}
\end{bt}

% nguồn: Ok_TuanLuan_Toan11_Lan1
\begin{bt}%[1C3Y1-2]%[ma:71bf2559]
	Tìm giới hạn sau $\lim\dfrac{2n^3-2n+3}{1-4n^3}$. 
	\loigiai{ 
	\[\lim\dfrac{2n^3-2n+3}{1-4n^3}=\lim\dfrac{2-\dfrac{2}{n^2}+\dfrac{3}{n^3}}{\dfrac{1}{n^3}-4}=-\dfrac{1}{2}.\]} 
\end{bt}

% nguồn: Ok_TuanLuan_Toan11_Lan1
\begin{bt}%[1C3Y1-2]%[ma:3ffb488c]
	Tìm  giới hạn sau $\lim\dfrac{\sqrt{n^4+2n+2}}{n^2+1}$.
	 \loigiai{
	\[\lim\dfrac{\sqrt{n^4+2n+2}}{n^2+1}  = \lim\dfrac{\sqrt{1+\dfrac{2}{n^3}+\dfrac{2}{n^4}}}{1+\dfrac{1}{n^2}}=1.\]}
\end{bt}

% nguồn: Ok_TuanLuan_Toan11_Lan1
\begin{bt}%[1C3B1-2]%[ma:7a7b32e0]
	 Tìm giới hạn sau $\lim\dfrac{1+2+2^2+\cdots +2^n}{1+3+3^2+\cdots +3^n}$.
	\loigiai{ 
	\[\lim\dfrac{1+2+2^2+\cdots +2^n}{1+3+3^2+\cdots +3^n} = \lim\dfrac{\dfrac{1-2^{n+1}}{-1}}{\dfrac{1-3^{n+1}}{-2}} = \lim\dfrac{\left(1-2^{n+1}\right)\cdot 2}{1-3^{n+1}} =  \lim\dfrac{\left(\left(\dfrac{1}{3}\right)^{n+1}-\left(\dfrac{2}{3}\right)^{n+1}\right)\cdot 2}{\left(\dfrac{1}{3}\right)^{n+1}-1}=0.\]}
\end{bt}

% nguồn: Ok_TuanLuan_Toan11_Lan1
\begin{bt}%[1C3B1-2]%[ma:f42e8620]
	Cho hai dãy số $(u_n), (v_n)$ với $u_n=3+\dfrac{1}{n}; v_n=5-\dfrac{2}{n^2}$.  Tính các giới hạn sau:
	\begin{enumEX}{1}
		\item $\lim u_n, \lim  v_n$.
		\item $\lim \left(u_n+v_n\right)$, $\lim  \left(u_n-v_n\right)$, $\lim \left(u_n\cdot v_n\right)$, $\lim \dfrac{u_n}{v_n}$.
	\end{enumEX}
	\loigiai{ 
		\begin{enumerate}
			\item Ta có  \\
			$\lim u_n=\lim \left(3+\dfrac{1}{n}\right)=\lim  3+\lim \left(\dfrac{1}{n}\right)=3+0=3$.\\
			$\lim v_n=\lim \left(5-\dfrac{2}{n^2}\right)=\lim  5-\lim \left(\dfrac{2}{n^2}\right)=5-0=5$.\\
			\item Ta có  \\
			$\lim \left(u_n+v_n\right)=\lim  u_n+\lim v_n=3+5=8$.\\
			$\lim \left(u_n-v_n\right)=\lim u_n-\lim v_n=3-5=-2$.\\ 
			$\lim \left(u_n\cdot v_n\right)=\lim  u_n\cdot \lim v_n=3\cdot 5=15$.\\
			 $\lim \dfrac{u_n}{v_n}=\dfrac{3}{5}$.
		\end{enumerate}
	} 
\end{bt}

% nguồn: Ok_TuanLuan_Toan11_Lan1
\begin{bt}%[1C3K1-2]%[ma:ea7e0b9f]
	Tính các giới  hạn sau:
	\begin{enumEX}{3}
		\item $\lim \dfrac{5n+1}{2n}$; 
		 \item $\lim \dfrac{6n^2+8n+1}{5n^2+3}$;
		\item $\lim \dfrac{\sqrt{n^2+5n+3}}{6n+2}$; 
		\item $\lim \left(2-\dfrac{1}{3^n}\right)$; 
		\item  $\lim \dfrac{3^n+2^n}{4\cdot3^n}$;
		\item  $\lim \dfrac{2+\dfrac{1}{n}}{3^n}$.
	\end{enumEX}
	 \loigiai{
		\begin{enumerate}
			\item Ta có $\lim \dfrac{5n+1}{2n}=\lim \dfrac{5+\dfrac{1}{n}}{2}=\dfrac{5}{2}$. 
			\item Ta có $\lim  \dfrac{6n^2+8n+1}{5n^2+3}=\lim\dfrac{6+\dfrac{8}{n}+\dfrac{1}{n^2}}{5+\dfrac{3}{n^2}}=\dfrac{6}{5}$.
			\item Ta có $\lim \dfrac{\sqrt{n^2+5n+3}}{6n+2}=\lim \dfrac{\sqrt{1+\dfrac{5}{n}+\dfrac{3}{n^2}}}{6+\dfrac{2}{n}}=\dfrac{1}{6}$. 
			\item Ta có $\lim \left(2-\dfrac{1}{3^n}\right)=\lim 2 -\lim \left(\dfrac{1}{3^n}\right)=2-0=2$. 
			\item  Ta có $\lim \dfrac{3^n+2^n}{4\cdot3^n}=\lim \dfrac{1+\left(\dfrac{2}{3}\right)^n}{4}=\dfrac{1}{4}$.
			\item Vì $\lim \left(2+\dfrac{1}{n}\right)=2$ và $\lim  3^n=+\infty$ nên $\lim \dfrac{2+\dfrac{1}{n}}{3^n}=0$.
		\end{enumerate}
	} 
\end{bt}

% nguồn: Ok_TuanLuan_Toan11_Lan2
\begin{ex}%[1C3Y1-2]%[ma:3228495f]
	Tìm giới hạn  sau $\lim\dfrac{3^{n+1}-4^n}{4^{n-1}+3}$.
	\loigiai{ 
	 \[\lim\dfrac{3^{n+1}-4^n}{4^{n-1}+3} = \lim\dfrac{9\cdot 3^{n-1}-4\cdot 4^{n-1}}{4^{n-1}+3}=\lim\dfrac{9\cdot\left(\dfrac{3}{4}\right)^{n-1}-4}{1+3\cdot\left(\dfrac{1}{4}\right)^{n-1}}=-4.\]}
\end{ex}

% nguồn: Ok_TuanLuan_Toan11_Lan2
\begin{ex}%[1C3Y1-2]%[ma:0ab9a0ca]
	 Tìm giới hạn sau $\lim\dfrac{2n^3-2n+3}{1-4n^3}$.
	 \loigiai{
	\[\lim\dfrac{2n^3-2n+3}{1-4n^3}=\lim\dfrac{2-\dfrac{2}{n^2}+\dfrac{3}{n^3}}{\dfrac{1}{n^3}-4}=-\dfrac{1}{2}.\]} 
\end{ex}

% nguồn: Ok_TuanLuan_Toan11_Lan2
\begin{ex}%[1C3Y1-2]%[ma:4ff19d6d]
	Tìm giới hạn  sau $\lim\dfrac{\sqrt{n^4+2n+2}}{n^2+1}$.
	\loigiai{ 
	 \[\lim\dfrac{\sqrt{n^4+2n+2}}{n^2+1} = \lim\dfrac{\sqrt{1+\dfrac{2}{n^3}+\dfrac{2}{n^4}}}{1+\dfrac{1}{n^2}}=1.\]}
\end{ex}

% nguồn: Ok_TuanLuan_Toan11_Lan2
\begin{ex}%[1C3B1-2]%[ma:a88b9650]
	Tìm giới hạn sau $\lim\dfrac{1+2+2^2+\cdots +2^n}{1+3+3^2+\cdots  +3^n}$.
	\loigiai{ 
	\[\lim\dfrac{1+2+2^2+\cdots +2^n}{1+3+3^2+\cdots +3^n} = \lim\dfrac{\dfrac{1-2^{n+1}}{-1}}{\dfrac{1-3^{n+1}}{-2}}  = \lim\dfrac{\left(1-2^{n+1}\right)\cdot 2}{1-3^{n+1}} = \lim\dfrac{\left(\left(\dfrac{1}{3}\right)^{n+1}-\left(\dfrac{2}{3}\right)^{n+1}\right)\cdot 2}{\left(\dfrac{1}{3}\right)^{n+1}-1}=0.\]}
\end{ex}

% nguồn: Ok_TuanLuan_Toan11_Lan2
\begin{ex}%[1C3B1-2]%[ma:8375c5ed]
	Cho hai dãy số $(u_n), (v_n)$ với $u_n=3+\dfrac{1}{n}; v_n=5-\dfrac{2}{n^2}$. Tính các  giới hạn sau:
	\begin{enumEX}{1}
		\item  $\lim u_n, \lim v_n$.
		\item $\lim \left(u_n+v_n\right)$, $\lim \left(u_n-v_n\right)$, $\lim \left(u_n\cdot v_n\right)$,  $\lim \dfrac{u_n}{v_n}$.
	\end{enumEX}
	 \loigiai{
		\begin{enumerate}
			\item Ta có  \\
			$\lim  u_n=\lim \left(3+\dfrac{1}{n}\right)=\lim 3+\lim \left(\dfrac{1}{n}\right)=3+0=3$.\\
			$\lim  v_n=\lim \left(5-\dfrac{2}{n^2}\right)=\lim 5-\lim \left(\dfrac{2}{n^2}\right)=5-0=5$.\\
			\item  Ta có \\
			$\lim \left(u_n+v_n\right)=\lim u_n+\lim  v_n=3+5=8$.\\
			$\lim  \left(u_n-v_n\right)=\lim u_n-\lim v_n=3-5=-2$.\\
			$\lim \left(u_n\cdot v_n\right)=\lim  u_n\cdot \lim v_n=3\cdot 5=15$.\\
			 $\lim \dfrac{u_n}{v_n}=\dfrac{3}{5}$.
		\end{enumerate}
	} 
\end{ex}

% nguồn: Ok_TuanLuan_Toan11_Lan2
\begin{ex}%[1C3K1-2]%[ma:c132f381]
	 Tính các giới hạn sau:
	\begin{enumEX}{3}
		\item $\lim \dfrac{5n+1}{2n}$; 
		 \item $\lim \dfrac{6n^2+8n+1}{5n^2+3}$;
		\item  $\lim \dfrac{\sqrt{n^2+5n+3}}{6n+2}$;
		 \item $\lim \left(2-\dfrac{1}{3^n}\right)$;
		\item  $\lim \dfrac{3^n+2^n}{4\cdot3^n}$;
		\item $\lim \dfrac{2+\dfrac{1}{n}}{3^n}$. 
	\end{enumEX}
	\loigiai{ 
		\begin{enumerate}
			\item Ta có $\lim  \dfrac{5n+1}{2n}=\lim \dfrac{5+\dfrac{1}{n}}{2}=\dfrac{5}{2}$.
			 \item Ta có $\lim \dfrac{6n^2+8n+1}{5n^2+3}=\lim\dfrac{6+\dfrac{8}{n}+\dfrac{1}{n^2}}{5+\dfrac{3}{n^2}}=\dfrac{6}{5}$.
			 \item Ta có $\lim \dfrac{\sqrt{n^2+5n+3}}{6n+2}=\lim \dfrac{\sqrt{1+\dfrac{5}{n}+\dfrac{3}{n^2}}}{6+\dfrac{2}{n}}=\dfrac{1}{6}$.
			\item Ta có  $\lim \left(2-\dfrac{1}{3^n}\right)=\lim 2 -\lim \left(\dfrac{1}{3^n}\right)=2-0=2$.
			\item Ta có $\lim \dfrac{3^n+2^n}{4\cdot3^n}=\lim  \dfrac{1+\left(\dfrac{2}{3}\right)^n}{4}=\dfrac{1}{4}$.
			\item Vì $\lim  \left(2+\dfrac{1}{n}\right)=2$ và $\lim 3^n=+\infty$ nên $\lim \dfrac{2+\dfrac{1}{n}}{3^n}=0$.
		\end{enumerate}
	 }
\end{ex}
